\section{Rappel : fonctions et blocs fonctionnels}
\label{secRappelBlocsFonctionnels}

Les chaînes de conversion étudiées dans ce document reposent sur des
\textbf{blocs fonctionnels} (\emph{function blocks}, norme IEC~61131-3).
Avant d'aller plus loin, il est utile d'en rappeler la nature exacte et de
la comparer à celle d'une simple fonction.

\begin{UPSTIinfor}{Fonction ou bloc fonctionnel ?}
\begin{description}
    \item[Fonction] Ne possède \textbf{aucune mémoire} entre deux appels :
    son résultat ne dépend que des arguments fournis à l'appel en cours.
    Elle ne peut retourner qu'une seule valeur (éventuellement composée),
    directement, par son nom.
    \item[Bloc fonctionnel] Est \textbf{instancié} (une variable, de type
    bloc fonctionnel, doit être déclarée) et \textbf{conserve un état}
    entre deux appels successifs : fronts montants mémorisés, étapes d'une
    temporisation, dernière erreur rencontrée... Il peut exposer plusieurs
    sorties simultanément (\texttt{VAR\_OUTPUT}) et, pour les blocs étudiés
    ici, des paramètres \texttt{VAR\_IN\_OUT} qu'il lit \emph{et} modifie.
\end{description}
\end{UPSTIinfor}

Tous les blocs rencontrés dans ce document et dans les TP suivants sont
des blocs fonctionnels. C'est en particulier le cas des blocs \texttt{MC\_}
(un bloc de mouvement doit savoir, d'un cycle à l'autre, s'il a déjà
démarré le mouvement demandé, afin de ne déclencher l'exécution que sur un
\textbf{front montant} de son entrée \texttt{Execute}) : une fonction,
sans mémoire, ne pourrait pas détecter ce front.

\begin{UPSTIactivite}[4][Fonction ou bloc fonctionnel ?][][][TD]
\UPSTIquestion{Les éléments suivants sont tous utilisés dans les TP de la
cellule. Pour chacun, indiquer s'il s'agit d'une \textbf{fonction} (appel
direct, résultat récupéré dans une variable) ou d'un \textbf{bloc
fonctionnel} (instance à réserver dans l'éditeur de données, onglet
\emph{Bloc fonction}), en justifiant en une phrase par le critère de
mémorisation d'état :
\begin{enumerate}
    \item \texttt{InitChart}, appelé au premier cycle pour initialiser le
    GEMMA, dont la valeur de retour est récupérée dans un objet
    \texttt{xResult : BOOL} ;
    \item \texttt{MC\_GroupReset}, déclenché sur front montant de son
    entrée \texttt{Execute} ;
    \item \texttt{TON}, temporisateur standard IEC~61131-3, dont la sortie
    \texttt{Q} passe à \texttt{TRUE} lorsque son entrée \texttt{IN} est restée
    à \texttt{TRUE} pendant au moins la durée \texttt{PT} ;
    \item la détection du front montant du bouton « Arrêt demandé » du
    pupitre.
\end{enumerate}}
\UPSTIquadrillage{4}

\begin{UPSTIcorrectionP}{Corrigé}
\begin{enumerate}
    \item \texttt{InitChart} est une \textbf{fonction} : elle est appelée
    directement (pas d'instance) et retourne une unique valeur
    \texttt{BOOL} récupérée dans \texttt{xResult} ; son action (rendre
    actif l'état initial du graphe passé en argument) ne nécessite aucune
    mémoire d'un appel à l'autre.
    \item \texttt{MC\_GroupReset} est un \textbf{bloc fonctionnel} : pour
    ne réagir qu'au \textbf{front montant} d'\texttt{Execute}, il doit
    mémoriser la valeur d'\texttt{Execute} au cycle précédent (et l'état
    d'avancement de sa requête) -- d'où l'instance
    \texttt{InstGroupResetD2} à réserver.
    \item \texttt{TON} est un \textbf{bloc fonctionnel} : pour savoir si
    \texttt{IN} est à \texttt{TRUE} depuis au moins \texttt{PT}, il doit
    mémoriser d'un cycle à l'autre le temps écoulé depuis lequel
    \texttt{IN} est à \texttt{TRUE} -- une fonction, qui ne voit que la
    valeur courante de \texttt{IN}, en serait incapable.
    \item La détection de front montant est un \textbf{bloc fonctionnel}
    (\texttt{R\_TRIG}) : détecter un front, c'est comparer la valeur
    courante à celle du cycle précédent, qui doit donc être mémorisée.
\end{enumerate}
\end{UPSTIcorrectionP}

\UPSTIquestion{Le mode D2 du TP~3 utilise les blocs
\texttt{MC\_GroupReset}, \texttt{MC\_GroupStop} et
\texttt{MC\_GroupContinue}. En respectant la règle de nommage imposée
(nom d'instance faisant référence au mode, par exemple
\texttt{InstGroupResetD2}), écrire la déclaration des trois instances
nécessaires.}

\begin{lstlisting}
VAR
    (*@\trou[InstGroupResetD2]{4cm}@*) : (*@\trou[MC\_GroupReset]{3.5cm}@*);
    (*@\trou[InstGroupStopD2]{4cm}@*) : (*@\trou[MC\_GroupStop]{3.5cm}@*);
    (*@\trou[InstGroupContinueD2]{4cm}@*) : (*@\trou[MC\_GroupContinue]{3.5cm}@*);
END_VAR
\end{lstlisting}
\end{UPSTIactivite}

\UPSTIremarque[Préfixes de bibliothèque]{Les blocs fonctionnels utilisés
sur la cellule suivent une convention de préfixe selon leur origine :
\texttt{VAL\_} pour la bibliothèque \texttt{uniVALplc} du robot,
\texttt{MC\_} pour les blocs \texttt{Motion Control} conformes au standard
PLCopen, \texttt{FB\_} pour les blocs maison de l'application.}
